%LTeX: language=es
\section{Ecuación de advección}

Volvemos sobre la ecuación de advección a velocidad constante \eqref{eq:adveccion}. Para más simplicidad supongamos que $d=1$ y $\bm v(t,x) = a$. Legamos a la ecuación
\begin{equation}
    \tag{A}
    \label{eq:adveccion 1d}
    \begin{dcases}
        \frac{\partial u}{\partial t}(t,x) + a \frac{\partial u}{\partial x}(t,x) = 0
         & \text{para todo }t > 0, x \in \R \\
        u(0,x) = u_0(x)
         & \text{para todo } x \in \R.
    \end{dcases}
\end{equation}
Usando desarrollos de Taylor de primer orden tenemos cuatro formas de discretizar la ecuación
\begin{example}[Algunos esquemas numéricos para \eqref{eq:adveccion}]
    \begin{align}
        \label{eq:adveccion 1d EE forward}
        \tag{A-EE$_+$}
         & \frac{U_i^{n+1} - U_i^n}{\tau} + a \frac{U_{i+1}^{n} - U_i^n}{h} = 0
        \\
        \tag{A-EE$_-$}
        \label{eq:adveccion 1d EE backward}
         & \frac{U_i^{n+1} - U_i^n}{\tau} + a \frac{U_{i}^{n} - U_{i-1}^n}{h} = 0
        \\
        \tag{A-EI$_+$}
         & \frac{U_i^{n+1} - U_i^n}{\tau} + a \frac{U_{i+1}^{n+1} - U_i^{n+1}}{h} = 0
        \\
        \tag{A-EI$_-$}
         & \frac{U_i^{n+1} - U_i^n}{\tau} + a \frac{U_{i}^{n+1} - U_{i-1}^{n+1}}{h} = 0.
    \end{align}
\end{example}
Las dos primeras discretizaciones son explícitas, y las dos segundas son implícitas. Vamos a centrarnos primero en las explícitas.


Observamos algunas propiedades de la solución exacta
\begin{enumerate}
    \item $\sup_x u(\cdot,t) = \sup u_0$ y lo mismo para el ínfimo. En particular, si $u_0 \ge 0$ entonces $u \ge 0$.
    \item Se preserva la integral porque
          \begin{equation*}
              \int_\R u(t,x) dx = \int_\R u_0 (t-ax) dx = \int_\R u_0(y) dy.
          \end{equation*}
\end{enumerate}

Vamos a estudiar el caso $a > 0$, y el otro se sigue por simetría.

La solución explícita claramente cumple . Veamos qué pasa con los diferentes métodos
\paragraph{Esquema \eqref{eq:adveccion 1d EE forward}.} Escribimos
\begin{align*}
    U_i^{n+1}
     & = U_i^n - a \tfrac{\tau}{h} (U_{i+1}^n - U_i^n)               \\
     & = (1 + a\tfrac{\tau}{h}) U_i^n - a \tfrac{\tau}{h} U_{i+1}^n.
\end{align*}
Si $a > 0$, $U_1^0 = 0$ y $U_2^0 = 1$, entonces $U_1^1 < 0$. Este método no preserva propiedades básicas.

\paragraph{Esquema \eqref{eq:adveccion 1d EE backward}.}
Escribimos
\begin{align*}
    U_i^{n+1}
     & = U_i^n - a \tfrac{\tau}{h} (U_{i}^n - U_{i-1}^n)             \\
     & = (1 - a\tfrac{\tau}{h}) U_i^n + a \tfrac{\tau}{h} U_{i-1}^n.
\end{align*}
Si $a > 0$ y $U_\bullet^0 \ge 0$ entonces el segundo término no es problemático. El primer término preserva la positividad si
\begin{equation}
    \tag{A-CFL}
    \label{eq:adveccion 1d CFL}
    a\tau \le h.
\end{equation}
Con esta condición tan inocente, resulta que el supremo está controlado, pues
\begin{equation*}
    \sup_i U_i^{n+1} \le (1 - a \tfrac{\tau}{h}) \sup_i U_i^n + a \tfrac{\tau}{h} \sup_i U_{i-1}^n = \sup_i U_i^{n}.
\end{equation*}
Lo mismo ocurre con el ínfimo. Es más, esto mismo se puede hacer con valores absolutos
\begin{equation*}
    \sup_i |U_i^{n+1}| \le \sup_i |U_i^{n}|.
\end{equation*}
Como el problema es lineal, si tenemos dos soluciones $U$ y $V$ tenemos que
\begin{equation*}
    \sup_i |U_i^{n} - V_i^n| \le \sup_i |U_i^0 - V_i^0|.
\end{equation*}

\paragraph{Soluciones aproximadas de \eqref{eq:adveccion 1d EE backward}.}
Si $V$ no satisface exactamente la ecuación, si no que tiene un ``resto'' $R$ en el sentido que
\begin{equation}
    \label{eq:adveccion 1d EE V con F}
    \frac{V_i^{n+1} - V_i^n}{\tau} + a \frac{V_{i}^{n} - V_{i-1}^n}{h} = R_i^n
\end{equation}
entonces la cuenta nos dice que
\begin{equation*}
    \sup_i |U_i^{n+1} - V_i^{n+1}| \le \sup_i |U_i^n - V_i^n| + \tau \sup_i |R_i^{n}|.
\end{equation*}
Iterando deducimos que
\begin{equation}
    \tag{A-S}
    \label{eq:adveccion 1d dependencia continua}
    \sup_i |U_i^{n+1} - V_i^{n+1}| \le \sup_i |U_i^0 - V_i^0| + \tau \sum_{k=0}^n \sup_i |R_i^{k}|.
\end{equation}
Esta ecuación nos dice que si tomamos

\paragraph{Consistencia de \eqref{eq:adveccion 1d EE backward}}

Definimos $V_i^n = u(t_n, x_i) = u_0(t_n - ax_i)$.
Tenemos \eqref{eq:adveccion 1d EE V con F} con
\begin{align*}
    R_i^{n} & \defeq \frac{V_i^{n+1} - V_i^n}{\tau} + a \frac{V_i^{n} - V_{i-1}^n}{\tau}
    \\
            & = \frac{V_i^{n+1} - V_i^n}{\tau} + a \frac{V_i^{n} - V_{i-1}^n}{\tau} - \left( \frac{\partial u}{\partial t}(t_n, x_i) + a\frac{\partial u}{\partial x}(t_n, x_i)   \right)
    \\
            & = \frac{V_i^{n+1} - V_i^n}{\tau} - \frac{\partial u}{\partial t}(t_n, x_i) + a \left( \frac{V_i^{n} - V_{i-1}^n}{\tau} - \frac{\partial u}{\partial x}(t_n, x_i) \right).
\end{align*}
Si $u_0 \in C^2_c$ entonces
\begin{align*}
    \left|\frac{V_i^{n+1} - V_i^n}{\tau} - \frac{\partial u}{\partial t}(t_n, x_i) \right|
     & \le
    \frac 1 2 \left|\frac{\partial^2 u }{\partial t^2} (\xi, x)\right| \tau \\
    \left|\frac{V_i^{n} - V_{i-1}^n}{\tau} - \frac{\partial u}{\partial x}(t_n, x_i) \right|
     & \le
    \frac 1 2 \left|\frac{\partial^2 u }{\partial x^2} (t_n, \eta)\right| h .
\end{align*}
Es decir que $|F_{i}^n| \le C(\tau + h)$. Aplicando \eqref{eq:adveccion 1d CFL} tenemos
\begin{equation}
    \tag{A-C}
    \label{eq:adveccion 1d consistencia}
    |R_i^n| \le C h.
\end{equation}

\paragraph{Convergencia de \eqref{eq:adveccion 1d EE backward}}

Así, aplicando \eqref{eq:adveccion 1d dependencia continua} tenemos, dado que $U_i^0 = u_0(x_i) = u(0, x_i) = V_i^0$ tenemos que sobre los puntos de la malla la solución numérica se parece a la exacta.
\begin{equation}
    \tag{A-E}
    \label{eq:adveccion 1d convergencia}
    \sup_i |U_i^{n+1} - u(t_{n+1}, x_i)| \le  C \tau n h.
\end{equation}

Nótese que $\tau n = t_n$. Hemos demostrado el siguiente teorema
\begin{teorema}
    Si $a > 0$, $u_0 \in C^2_c(\mathbb R)$, y supongamos \eqref{eq:adveccion 1d CFL}. Entonces la solución de \eqref{eq:adveccion 1d EE backward} con $U_i^0 = u_0 (ih)$ aproxima uniformemente la solución del problema \eqref{eq:adveccion 1d} con estimación \eqref{eq:adveccion 1d convergencia}.
\end{teorema}




